\documentclass[legalpaper,11pt]{article}
\usepackage[spanish, mexico, activeacute]{babel}
\usepackage{amsmath,amsfonts,amssymb}
\usepackage{fancyhdr}
\usepackage{fancyvrb}
\usepackage{xy}
\usepackage{graphicx}
\usepackage{latexsym}
\usepackage{enumerate}
\usepackage{tabularx}
\usepackage{multirow}
\usepackage{multicol}
\usepackage{hyperref}

\setlength{\oddsidemargin}{-0.5cm}
\setlength{\evensidemargin}{0cm}
\setlength{\textwidth}{17.5cm}
\setlength{\textheight}{31cm}
\setlength{\topmargin}{-2.5cm}

\font\bff=cmbx10 at 10truept
\font\lg=cmdunh10 at 10truept
\font\bl=cmss10 at 10truept

\newcommand{\matlab}{{\sc Matlab} }

\newcommand{\header}{\noindent%
\begin{minipage}{0.09\textwidth}
\includegraphics[width=\textwidth]{logoudec.eps}
\end{minipage}
\hspace{3mm}
\begin{minipage}{0.5\textwidth}
{\lg UNIVERSIDAD DE CONCEPCI\'ON}\hfill
\vskip-4truept
\noindent{\bff FACULTAD DE CIENCIAS}\hfill
\vskip-4truept
\noindent{\bff F\'ISICAS Y MATEM\'ATICAS}\hfill
\vskip-4truept
\noindent{\bl DEPARTAMENTO DE INGENIER\'IA MATEM\'ATICA}\hfill
\vskip4truept\hrule\hrule\vskip4truept
\par
\end{minipage}
}

\newcommand{\respuesta}[1]{
%\noindent\makebox[\textwidth][r]{
\fbox{
\begin{minipage}{\linewidth}
\hfill
\vspace{#1}
\end{minipage}
}
%}
}

\newcommand{\puntaje}[1]{\textbf{(#1 \%)}}

\newcommand{\titulo}[2]{
\newpage
\header
\begin{minipage}{0.4\textwidth}
	\begin{center}
    \ifnum#2=2
       \begin{tabular}{||p{2cm}|p{2.1cm}||}
       \hline
       Pregunta a &  Pregunta b   \\
       \hline
       \vspace{1cm}\null &     \\
       \hline
       \end{tabular}
    \else
    	\ifnum#2=3
       \begin{tabular}{||p{1.8cm}|p{1.8cm}|p{1.8cm}||}
       \hline
       Preg. a &  Preg. b &  Preg. c   \\
       \hline
       \vspace{1cm}\null &    & \\
       \hline
       \end{tabular}
     \else
       \ifnum#2=4
       \begin{tabular}{||p{1cm}|p{1cm}|p{1cm}|p{1cm}||}
       \hline
       Preg. a &  Preg. b &  Preg. c & Preg. d  \\
       \hline
       \vspace{1cm}\null &    & & \\
       \hline
       \end{tabular}
    \else
       \ifnum#2=5
       \begin{tabular}{||p{0.8cm}|p{0.8cm}|p{0.8cm}|p{0.8cm}|p{0.8cm}||}
       \hline
       Preg. a &  Preg. b &  Preg. c & Preg. d &Preg. e \\
       \hline
       \vspace{1cm}\null &    & & &\\
       \hline
       \end{tabular}
     \else
       \begin{tabular}{||p{2cm}||}
       \hline
       Pregunta a   \\
       \hline
       \vspace{1cm}\null      \\
       \hline
       \end{tabular}
    \fi
    \fi
    \fi
    \fi
  \end{center}
\end{minipage}
\begin{center}
\textbf{\Large TEST 2 VERSI\'ON #1 C\'ALCULO NUM\'ERICO 521230}\\
\vspace{5mm}
 \begin{tabular}{p{0.6\textwidth}p{0.35\textwidth}}
	\textbf{Nombre:}    &\textbf{Carrera:}\\ \\
	\textbf{Ayudante:} & \textbf{Matr\'icula:}
 \end{tabular}
 \\
\end{center}
 
 Enviar los programas solicitados en el formato solicitado al correo informado por el ayudante de su secci\'on y con copia a \textbf{numerico@ing-mat.udec.cl}.
 }
 
\usepackage[numbered,framed]{matlab-prettifier}
\lstMakeShortInline"
\lstset{
  style              = Matlab-editor,
  escapechar         = ",
  mlshowsectionrules = true,
}

 
 \pagestyle{empty}


\title{T2 FINAL 521230 2018-1}
\begin{document}
\titulo{1}{5}	%Versión,cantidad de preguntas 				
\begin{enumerate}
\item 
 Considere el P.V.I.
$$
\begin{array}{rl|}
xy'''+2xy'	&=0 \\
y(0)	&=1\\
y'(0)	&=1\\
y''(0)	&=2\\
\hline
\end{array}
$$
cuya soluci\'on exacta es 
$$
y(x) = \frac{sen(\sqrt{2} x)}{\sqrt{2}} - cos(\sqrt{2} x) + 2.
$$
En el siguiente casillero, mediante sustituciones, reduzca el orden del P.V.I. para ser resuelto usando los m\'etodos RK.

\hspace{-10mm}
\framebox(500,250){}

En un programa de Matlab llamado \texttt{programa1.m}
\begin{enumerate}
\item\puntaje{30}  Resuelva num\'ericamente el P.V.I. usando el m\'etodo de Euler (RK11) y el m\'etodo RK45 usado en \texttt{ode45} sobre la malla \texttt{0:0.1:10}.
\item \puntaje{20}  Grafique en un mismo gr\'afico la soluci\'on exacta y las obtenidas anteriormente.
\item\puntaje{20}  Calcule los errores locales de estos m\'etodos.
\item  \puntaje{15} Grafique los pares ordenados $\{(x_i,\tau_i)\}_{i=1}^N$ obtenidos, donde $x_i$ es el $i$-\'esimo nodo y $\tau_i$ es el $i$-\'esimo error local para cada soluci\'on num\'erica calculada.
\item \puntaje{15} Determine el error global de cada m\'etodo.
\end{enumerate}
Adjunte el programa al correo.

\textbf{Desarrollo:} La sustituci\'on necesaria para este problema es
$$
u_1=y, \quad u_2=y', \quad u_3=y'',
$$
con la cual se crea el sistema
$$
\begin{array}{cl|}
u_1'	&=u_2\\
u_2'	&=u_3\\
u_3'	&=-2 u_2\\ 
\hline
\end{array}
\text{ con las condiciones iniciales }
\begin{array}{cl|}
u_1(0)	&=1\\
u_2(0)	&=1\\
u_3(0)	&=2\\ 
\hline
\end{array}
$$
El programa solicitado debe tener instrucciones similares a
\begin{lstlisting}
%%Método de Euler
DF=@(x,u) [0,1,0;0,0,1;0,-2,0]*u
x=0:0.1:1;
u=[1;1;2]
for i=1:length(x)-1
   u(:,i+1)=DF(x(i),u(:,i))*(x(i+1)-x(i))+u(:,i);
endfor

%Método de ode45
[X,U]=ode45(DF,x,[1;1;2]);    %%30 PUNTOS

plot(x,u(1,:),X,U(:,1))   	%%20 PUNTOS

f=@(x) sin(sqrt(2)*x)/sqrt(2)-cos(sqrt(2)*x)+2;
errorLocal1=abs(u(1,:)-f(x));
errorLocal2=abs(U(:,1)-f(X)); %%20 PUNTOS

figure(2)
subplot(1,2,1)
plot(x,errorLocal1);
subplot(1,2,2)
plot(x,errorLocal2);  %%15 PUNTOS

errorGlobal1=max(errorLocal1)
errorGlobal2=max(errorLocal2) %%15 PUNTOS
\end{lstlisting}
con las cuales se hacen las gr\'aficas
\begin{center}
\includegraphics[width=0.6\textwidth]{./T2V1-1.eps}

\includegraphics[width=0.6\textwidth]{./T2V1-2.eps}
\end{center}



\newpage
\titulo{2}{2}
\item   Dado $N \in \mathbb{N} = \{1, 2, 3, \dotsc\}$ se define la siguiente regla de cuadratura sobre el intervalo $[0,2\pi]$:
%
\begin{subequations}\label{trig-quad}
\begin{equation}
\int_0^{2\pi} f(x) \, \mathrm{d}x \approx I_N(f) := \sum_{k=0}^{N-1} w_k \, f(x_k),
\end{equation}
%
donde
%
\begin{equation}
(\forall\,k\in\{0,\dotsc,N-1\}) \qquad x_k := \frac{2\pi k}{N} \quad\text{y}\quad w_k := \frac{2\pi}{N}.
\end{equation}
\end{subequations}
%
Note que el extremo $0$ es un nodo ($x_0)$ pero el extremo $2\pi$ no lo es.

Se dice la regla de cuadratura $I_N$ es exacta para polinomios trigonom\'etricos de grado menor o igual a $M \in \mathbb{N}$ si y solo si la regla integra exactamente a las funciones en la colecci\'on $\{1\} \cup \bigcup_{m=1}^M \{ \cos(mx), \sen(mx) \}$; esto es,
%
\begin{equation*}
I_N(1) = 2\pi \qquad\wedge\quad (\forall\,m\in\{1,\dotsc,M\}) \quad I_N(\cos(mx)) = 0 \ \wedge\ I_N(\sen(mx)) = 0.
\end{equation*}

\begin{enumerate}
\item \puntaje{60} Escriba una funci\'on \matlab que reciba una funci\'on $f$ y un $N \in \mathbb{N}$ y devuelva la cuadratura $I_N(f)$.

Por ejemplo, si su funci\'on se llama \texttt{trigquad}, deber\'ia poder usarse de la siguiente manera:
\begin{lstlisting}
>> trigquad(@(x) x.^0, 10) % @(x) x.^0 es la funci\'on 1 vectorizada

ans =

    6.2832

>> trigquad(@(x) cos(3*x), 5)

ans =

   1.3951e-16
\end{lstlisting}

\medskip

?`C\'omo llam\'o a su funci\'on?

\begin{tabularx}{\linewidth}{|p{7em}|X|}\hline
Nombre funci\'on & \\\hline
\end{tabularx}

\medskip

\item \puntaje{40} Escriba un rutero \matlab donde se compruebe num\'ericamente que la regla de cuadratura $I_N$ es exacta (salvo errores de redondeo) para polinomios trigonom\'etricos de grado menor o igual a $N-1$ pero \textbf{no lo es} para polinomios trigonom\'etricos de grado menor o igual a $N$, para $N \in \{1, 6, 11\}$.

\medskip

?`C\'omo llam\'o a su rutero?

\begin{tabularx}{\linewidth}{|p{7em}|X|}\hline
Nombre rutero & \\\hline
\end{tabularx}
\end{enumerate}

\textbf{Desarrollo:}
\begin{enumerate}
\item La funci\'on solicitada debe tener instrucciones similares a
\begin{lstlisting}
function val = trigquad (f,N)
  x=2*pi/N*[0:N-1];
  w=2*pi/N*ones(1,N);
  val=w*f(x)';
endfunction
\end{lstlisting}
\hfill\fbox{60 puntos}

\item El programa debe contener instrucciones similares a
\begin{lstlisting}
for N=[1,6,11]
    p=@(x) x.^0
    int=2*pi;
    quad=trigquad(p,N);
    if(abs(int-quad)>10^-10)
        disp(['Error de cuadratura :', num2str(abs(int-quad))] );
        disp(['Con nodos :', num2str(N)] );
        disp(['Para la función trigonometrica: f(x)=1']);
        endif
    for i=1:N
      p=@(x) cos(i*x);
      int=(sin(i*2*pi)-sin(i*0))/i;
      quad=trigquad(p,N);
      if(abs(int-quad)>10^-10)
        disp(['Error de cuadratura :', num2str(abs(int-quad))] );
        disp(['Con nodos :', num2str(N)] );
        disp(['Para la función trigonometrica: cos(',num2str(i),'x)']);
      endif
    endfor
    for i=1:N
      p=@(x) sin(i*x);
      int=(-cos(i*2*pi)+cos(i*0))/i;
      quad=trigquad(p,N);
      if(abs(int-quad)>10^-10)
        disp(['Error de cuadratura :', num2str(abs(int-quad))] );
        disp(['Con nodos :', num2str(N)] );
        disp(['Para la función trigonometrica: sen(',num2str(i),'x)']);
      endif    
    endfor
  endfor
\end{lstlisting}
\hfill\fbox{40 puntos}

El cual al ejecutarse muestra la salida

\begin{lstlisting}
Error de cuadratura :6.2832
Con nodos :1
Para la función trigonometrica: cos(1x)
Error de cuadratura :6.2832
Con nodos :6
Para la función trigonometrica: cos(6x)
Error de cuadratura :6.2832
Con nodos :11
Para la función trigonometrica: cos(11x)
\end{lstlisting}

\end{enumerate}


\titulo{3}{2}
\item Dado $N \in \mathbb{N} = \{1, 2, 3, \dotsc \}$ se define la siguiente regla de cuadratura de $N+1$ puntos sobre el intervalo $[-1,1]$:
%
\begin{subequations}\label{cheb-quad}
\begin{equation}\label{cheb-quad-main}
\int_{-1}^1 f(x) \, \frac{1}{\sqrt{1-x^2}} \, \mathrm{d}x
\approx I_N(f) := \sum_{k=0}^N w_k \, f(x_k),
\end{equation}
%
donde los nodos y los pesos est\'an dados por
%
\begin{equation}
(\forall\,k\in\{0,\dotsc,N\}) \quad x_k = \cos\left(\frac{k \pi}{N}\right)
\end{equation}
%
y
%
\begin{equation}
w_0 = \frac{\pi}{2N}, \qquad (\forall\,k\in\{1, \dotsc, N-1\}) \quad w_k = \frac{\pi}{N}, \qquad w_N = \frac{\pi}{2N}.
\end{equation}
\end{subequations}

\textbf{Es muy importante} notar que las integrales que se aproximan (lado izquierdo de \eqref{cheb-quad-main}) involucran la \emph{funci\'on peso} $\frac{1}{\sqrt{1-x^2}}$ pero que la regla de cuadratura precisa solamente de evaluaciones de la funci\'on que en el integrando multiplica a dicha funci\'on peso.

\begin{enumerate}
\item \puntaje{60} Escriba una funci\'on \matlab que reciba una funci\'on $f$ y un $N \in \mathbb{N}$ y devuelva la cuadratura $I_N(f)$ para aproximar la integral que aparece en \eqref{cheb-quad-main}.

Por ejemplo, si su funci\'on se llama \texttt{weightedquad}, deber\'ia poder usarse de la siguiente manera:
\begin{Verbatim}[frame=single]
>> weightedquad(@(x) x.^0, 10) % "@(x) x.^0" es la funci\'on 1 vectorizada

ans =

    3.1416

>> weightedquad(@(x) x.^3, 5)

ans =

   1.3951e-16
\end{Verbatim}

\medskip3

?`C\'omo llam\'o a su funci\'on?

\begin{tabularx}{\linewidth}{|p{7em}|X|}\hline
Nombre funci\'on & \\\hline
\end{tabularx}

\medskip

\item \puntaje{40} Escriba un rutero \matlab donde se compruebe num\'ericamente que la regla de cuadratura $I_N$ es exacta (salvo errores de redondeo) para computar la integral del producto de cualquier polinomio de grado menor o igual a $2N-1$ por la funci\'on peso $\frac{1}{\sqrt{1-x^2}}$ pero \textbf{no lo es} para el computar la integral del producto de polinomios de grado $2N$ y la misma funci\'on peso para $N \in \{1, 6, 11\}$.

\emph{Indicaci\'on:} Note que
%
\begin{equation*}
\int_{-1}^1 \frac{1}{\sqrt{1-x^2}} \, \mathrm{d}x = \pi
\end{equation*}
%
y que
%
\begin{equation*}
(\forall\,n\in\mathbb{N}) \quad
\int_{-1}^1 x^n \frac{1}{\sqrt{1-x^2}} \, \mathrm{d}x
= \begin{cases}
\frac{2\sqrt{\pi}}{n} \frac{\Gamma\left(\frac{n+1}{2}\right)}{\Gamma\left(\frac{n}{2}\right)} & \text{si $n$ es par,}\\
0 & \text{si $n$ es impar.}
\end{cases}
\end{equation*}
%
En \matlab la funci\'on $\Gamma$ se llama \texttt{gamma}.
En \matlab se puede testear si un n\'umero \texttt{n} es par mediante \texttt{n == 2*round(n/2)}.

\medskip

?`C\'omo llam\'o a su rutero?

\begin{tabularx}{\linewidth}{|p{7em}|X|}\hline
Nombre rutero & \\\hline
\end{tabularx}

\medskip

\end{enumerate}
\textbf{Desarrollo:} 
\begin{enumerate}
\item 
La funci\'on debe tener una estructura similar a
\begin{lstlisting}
function val = weightedquad (f,N)
x=cos(pi/N*[0:N]);
w=[pi/2/N,pi/N*ones(1,N-1),pi/2/N];
val=w*f(x)';
endfunction
\end{lstlisting}
\hfill\fbox{60 puntos}

\item El rutero debe hacer instrucciones similares a
\begin{lstlisting}
for N=[1,6,11]
    for i=1:2*N
      p=@(x) x.^i;
      if i==2*round(i/2)
        int=2*sqrt(pi)/i*gamma((i+1)/2)/gamma(i/2);
      else
        int=0;
      endif
      quad=weightedquad(p,N);
      if(abs(int-quad)>=10^-10)
        disp(['Error de cuadratura :', num2str(abs(int-quad))] );
        disp(['Con nodos :', num2str(N)] );
        disp(['Para la función trigonometrica: x^',num2str(i)]);
      endif
    endfor
  endfor
\end{lstlisting}
\hfill\fbox{40 puntos}

con los cuales se observa que
\begin{lstlisting}
Error de cuadratura :1.5708
Con nodos :1
Para la función trigonometrica: x^2
Error de cuadratura :0.001534
Con nodos :6
Para la función trigonometrica: x^12
Error de cuadratura :1.498e-006
Con nodos :11
Para la función trigonometrica: x^22
\end{lstlisting}
\end{enumerate}

\newpage
\titulo{4}{4}
\item El P.V.C.
$$
\begin{array}{cl|}
y''(t)	&=-0.5y'(t)-9.81\\
y(0)	&=1\\
y(10)	&=10\\ \hline
\end{array}
$$
modela la altura de un objeto $y(t)$ en el instante $t$. El cual parte en la posici\'on $(0,1)$ y termina en la posici\'on $(10,10)$. Recuerde que este P.V.C. se resuelve num\'ericamente resolviendo el sistema de ecuaciones tridiagonal
$$
\begin{bmatrix}
2 & -1+\frac{h}{4} & 0 & \cdots & 0 \\
-1-\frac{h}{4} & 2 & -1+\frac{h}{4} & \cdots& 0 \\
\vdots & \ddots & \ddots & \ddots &\vdots \\
0 & \cdots &-1-\frac{h}{4} & 2 & -1+\frac{h}{4} \\
0 & \cdots & 0 &-1-\frac{h}{4} & 2  \\
\end{bmatrix}
\begin{bmatrix}
y_1\\y_2\\\vdots \\y_{n-2}\\y_{n-1}
\end{bmatrix}
=
h^2\begin{bmatrix}
9.81-1+\frac{h}{4}\\
9.81\\
\vdots \\
9.81\\
9.81-1-\frac{h}{4}
\end{bmatrix}
$$
el cual da origen a la soluci\'on num\'erica $\{(x_i,y_i)\}_{i=0}^n$ mediante el m\'etodo de diferencias finitas, en una malla de $n+1$ nodos equiespaciados definidos por $x_i=1+i\frac{9}{n}$, $\forall i \in\{0,\cdots, n\}$, donde $h=\frac{9}{n}$ es el tama\~{n}o de la malla del m\'etodo.

\begin{enumerate}
\item \puntaje{40} Cree una funci\'on de llamada \texttt{matricesDF.m}, que tenga como entrada el m\'aximo \'indice de nodos \texttt{n} y que retorne la matriz de coeficientes del sistema anterior y el vector del lado derecho.

\item \puntaje{40} En un rutero llamado \texttt{DF.m}, usando la funci\'on \texttt{matricesDF} resuelva el problema anterior para $n\in \{5,10,15,20,25,30\}$

\item \puntaje{10} En este mismo rutero, grafique las soluciones obtenidas en un mismo gr\'afico.

\item \puntaje{10} Determine la altura m\'axima que alcanza este objeto en su trayectoria.

\hspace{-10mm}
\framebox(500,100){}
\end{enumerate}

Adjunte el rutero \texttt{DF.m} y la funci\'on \texttt{matricesDF.m} al correo.

\textbf{Desarrollo:} 

La funci\'on \texttt{matricesDF} debe tener instrucciones similares a
\begin{lstlisting}
function [M,b]=matricesDF(n)
  h=9/n;
  M=diag(2*ones(1,n-1))+diag((-1-h/4)*ones(1,n-2),-1)+diag((-1+h/4)*ones(1,n-2),1);
  b=9.81*ones(n-1,1);
  b(1)=b(1)-1+h/4;
  b(end)=b(end)-1-h/4;
  b=h^2*b;
\end{lstlisting}
\hfill\fbox{40 puntos}

y el rutero debe contener las instrucciones
\begin{lstlisting}
clear all; close all; clc;
for n=5:5:30
   [M,b]=matricesDF(n);
   y=M\b;
   y=[1;y;10];
   t=1:9/n:10;
   plot(t,y);hold on;
end

max(y)
\end{lstlisting}
\hfill\fbox{ 50 puntos}

con la cual se observa que la altura m\'axima del objeto es alrededor de \texttt{79} unidades de medida.  \hfill\fbox{ 10 puntos}
 
 
 
 \newpage
 \titulo{5}{3}
 \item Considere el siguiente Problema de Valores Iniciales:

$$
\left\{
\renewcommand*{\arraystretch}{1.5}
\begin{array}{l}
 x'=3\sen(t)x+2ty+1 \qquad t\in[0, 1{.}5]\\
y''=2x+t^2y'-5y+e^t\qquad t\in[0, 1{.}5]\\
x(0)=0,\quad y(0)=2,\quad y'(0)=0\\
\end{array}\right.
$$

\noindent donde $x=x(t)$ e $y=y(t)$. Para resolver este problema, utilizaremos el M\'etodo de Euler Impl\'icito, en cual considera una partici\'on del intervalo $[0, 1{.}5]$ en $N$ subintervalos de tama\~no $h$, donde:
$$t_i=(i-1)h,\,\,\,i=1,\ldots,N+1\qquad \textrm{con} \qquad h=\dfrac{1{.}5}{N}. $$
  \begin{enumerate}
 \item \puntaje{20} En el siguiente casillero, mediante sustituciones, reduzca el orden del PVI para ser resuelto usando el m\'etodo de Euler Impl\'icito.

\hspace{-10mm}
\framebox(500,200){}

\item \puntaje{20} En el siguiente casillero escriba, como un sistema de ecuaciones lineales, el esquema resultante al aplicar el m\'etodo de Euler Impl\'icito al PVI del item anterior.

\hspace{-10mm}
\framebox(500,150){}

\item  Escriba un rutero \matlab, llamado \texttt{version5.m}, que realice las siguientes tareas:
\begin{itemize}
\item \puntaje{40} resuelva num\'ericamente el problema usando el esquema de Euler Impl\'icito obtenido en el item anterior, 
\item \puntaje{20} grafique las soluciones $x$ e $y$, en un mismo gr\'afico,
\end{itemize}
Puede escribir los programas adicionales que considere necesario para desarrollar este ejercicio, no olvide adjuntarlos al correo.
 \end{enumerate}
 
 \textbf{Desarrollo:}
 \begin{enumerate}
\item  Mediante las sustituciones
 $$
 u_1=x, \quad u_2=y, \quad u_3=y',
 $$
 el P.V.I. puede ser planteado como
 $$
 \begin{array}{l|}
 u_1'=3sen(t)u_1+2tu_2+1\\
 u_2'=u_3\\
 u_3'=2u_1+t^2u_3-5u_2+e^t\\
 u_1(0)=0, \quad u_2(0)= 2, \quad u_3(0)=0\\ \hline 
 \end{array}
 $$
\hfill\fbox{20 puntos}

\item Usando el m\'etodo de Euler impl\'icito
$$
\begin{bmatrix}
u_1^{i+1}\\ u_2^{i+1}\\ u_3^{i+1}
\end{bmatrix}
=
\begin{bmatrix}
3sen(t_{i+1})u_1^{i+1}+2t_{i+1}u_2^{i+1}+1\\
u_3^{i+1}\\
2u_1^{i+1}+t_{i+1}^2u_3^{i+1}-5u_2^{i+1}+e^{t_{i+1}}\\
\end{bmatrix}
\begin{bmatrix}
h
\end{bmatrix}
+
\begin{bmatrix}
u_1^{i}\\ u_2^{i}\\ u_3^{i}
\end{bmatrix}
$$
de donde, reordenando se llega al sistema
$$
\begin{bmatrix}
-3sen(t_{i+1})h+1 	& 2t_{i+1}h	& 0 \\
0					& 1 		& -h\\
2h					& -5h		&-t_{i+1}^2h
\end{bmatrix}
\begin{bmatrix}
u_1^{i+1}\\u_2^{i+1}\\u_3^{i+1}
\end{bmatrix}
=
\begin{bmatrix}
u_1^{i}+h\\u_2^{i}\\u_3^{i}+e^{t_{i+1}}h
\end{bmatrix}
$$
\hfill\fbox{20 puntos}

\item De donde se puede programar este m\'etodo similar a
\begin{lstlisting}
clear all; close all;
N=10
h=1.5/N
t=0:h:1.5;

u(1,1)=0;
u(2,1)=2;
u(3,1)=0;

for i=1:length(t)-1
  u(:,i+1)=[-3*sin(t(i+1))*h+1,2*t(i+1)*h,0; 0,1,-h;2*h,-5*h,-t(i+1)^2*h]\[u(:,i)+[h;0;exp(t(i+1))*h]]
endfor				%%40 PUNTOS

plot(t,u(1,:),t,u(2,:))		%% 20 PUNTOS
\end{lstlisting}
 
\end{enumerate}
\end{enumerate}
\end{document}
